% --- SECCIÓN ------------------------------------------------------------------
% --- DEFINICIONES -------------------------------------------------------------
\chapter{Definiciones}

Definiciones extraídas de los libro \cite{Beardon1991}, \cite{Devaney2022} and
\cite{Devaney2020}.

\begin{definicion}[Plano complejo extendido]
  El \emph{plano complejo extendido} es la unión
  \begin{equation*}
    \ExtC = \C \cup \{\infty\},
  \end{equation*}
  donde $\infty$ es un punto \enquote{abstracto} fuera de $\C$.
\end{definicion}

Definamos ahora una métrica en $\ExtC$ que trate todos los puntos con igual
facilidad y haga que $\infty$ pierda cualquier significado especial: \emph{la
métrica cordal}.

Empecemos definiendo la \emph{esfera de Riemann}. Para ello, identificamos $\C$
con el plano horizontal
\begin{equation*}
  \{ (x_1, x_2, x_3) \in \R^3 \mid x_3 = 0 \}
\end{equation*}
en $\R^3$. Sea $S$ la esfera unidad en $\R^3$ centrada en el origen, y sea $N =
(0,0,1)$ el polo norte.

Proyectamos ahora cada punto de $\C$ sobre un punto de $S$ tomando la
intersección de la esfera unidad $S$ con la recta que pasa por $N$ y el punto
$(x_1, x_2, 0)$. En la Figura~\ref{fig:stereographic_projection} se muestra una
ilustración de esta proyección.

\begin{figure}[htbp]
  \centering
  \includegraphics[width=\textwidth]{beardons_book/stereographic_projection.png}
  \caption{Proyección estereográfica.}
  \label{fig:stereographic_projection}
\end{figure}

Identificamos esta proyección como:
\begin{equation*}
  \pi : \ExtC \longrightarrow S,
\end{equation*}
donde para $z = x + \I y \in \C$:
\begin{equation*}
  \pi(z) = \left( \frac{2x}{1+|z|^2}, \frac{2y}{1+|z|^2}, \frac{|z|^2-1}{1+|z|^2} \right),
\end{equation*}
y $\pi(\infty) = (0,0,1)$.

\begin{definicion}[Métrica cordal]
  Dada la proyección $\pi$ de $\ExtC$ en $S$, definimos la \emph{métrica
  cordal}, $\sigma : \ExtC \times \ExtC \implies \R^+ \cup \{0\}$, como la
  distancia euclídea de $\pi(z)$ a $\pi(w)$:
  \begin{equation*}
    (z,w) \mapsto \sigma(z,w) = |\pi(z) - \pi(w)|.
  \end{equation*}
  Explícitamente, cuando $z, w \in \C$:
  \begin{equation*}
    \sigma(z,w) = \frac{2|z-w|}{(1+|z|^2)^{1/2}(1+|w|^2)^{1/2}}.
  \end{equation*}
  Para el punto del infinito, tenemos:
  \begin{equation*}
    \sigma(z,\infty) = \lim_{w \to \infty} \sigma(z,w) = \frac{2}{(1+|z|^2)^{1/2}}.
  \end{equation*}
\end{definicion}

\begin{observacion}[$z \mapsto 1/z$ es una isometría]
  $\sigma(z,w) = \sigma(z^{-1}, w^{-1})$, es decir, la aplicación $z \mapsto
  \frac{1}{z}$ es una isometría.
\end{observacion}

\begin{observacion}[Rotación en $S$ $\mapsto$ isometría en $\ExtC$]
  Cualquier rotación $\varphi: S \implies S$ de la esfera induce una aplicación
  $\pi^{-1} \varphi \pi : \ExtC \implies \ExtC$ que es una isometría. Estas
  isometrías son transformaciones de Möbius de la forma:
  \begin{equation*}
    z \mapsto \frac{a z - \bar{c}}{c z + \bar{a}},
  \end{equation*}
  donde $|a|^2 + |c|^2 = 1$.
\end{observacion}

\begin{definicion}[Función Racional]
  Una \emph{función racional} es una función $R: \ExtC \to \ExtC$ de la forma
  \begin{equation*}
    R(z) = \frac{P(z)}{Q(z)},
  \end{equation*}
  donde $P(z)$ y $Q(z)$ son polinomios, sin ser ambos nulos.
\end{definicion}

Supongamos que ni $P$ ni $Q$ son el polinomio nulo, en este caso $R$ no está
definida en las raíces comunes de $P$ y $Q$. Sin embargo, si estas existen,
podemos cancelar los factores comunes de $P$ y $Q$, obteniendo así dos
polinomios que no tendrán raíces comunes $P_1$ y $Q_1$. De ahora en adelante
consideraremos que $P$ y $Q$ no tienen raíces comunes, y si las tuvieran,
consideraríamos $R = P_1 / Q_1$. En estas circunstancias seguiremos los
convenios habituales con el punto del infinito en la esfera de Riemann: si $Q$
es el polinomio nulo, entonces $R$ es la función constante $\infty$; si $Q(z) =
0$ y $P$ no es el polinomio nulo, entonces $R(z) = \infty$; y tomaremos
$R(\infty) = \lim_{z \to \infty} R(z)$.

También nos será útil definir el grado de $R$ como el máximo entre los grados
de los polinomios $P$ y $Q$, con $P$ y $Q$ sin raíces comunes. Si $R$ es
constante, $R = \alpha \in \ExtC$, consideraremos que tiene grado~$0$.

Denotaremos la n-ésima composición de una función $f$ consigo misma en un punto
$z$ con~$f^n(z)$.

\begin{definicion}[Punto fijo y periódico]
  Dada una función $f$ y un punto $\zeta$, $\zeta$ es un \emph{punto fijo} para
  $f$ si $f(\zeta) = \zeta$. El punto $\zeta$ es un \emph{punto periódico} de
  periodo $n$ si $f^n(\zeta) = \zeta$. El menor $n$ positivo tal que
  $f^n(\zeta) = \zeta$ se llama \emph{periodo primo} de $\zeta$. Denotaremos el
  conjunto de puntos periódicos de periodo $n$ (no necesariamente primo) por
  $\periodicpoints{n}(f)$, y el conjunto de puntos fijos por $\fixedpoints(f)$.
  El conjunto de todas las iteraciones de un punto periódico es una
  \emph{órbita periódica}.
\end{definicion}

\begin{definicion}[Punto eventualmente periódico]
  Un punto $z$ es \emph{eventualmente periódico} de periodo $n$ si $z$ no es
  periódico pero existe $m > 0$ tal que $f^{i+n}(z) = f^i(z)$ para todo $i \geq
  m$. Es decir, $f^i(z)$ es periódico si $i \geq m$.
\end{definicion}

\begin{definicion}[Punto asintótico y conjunto estable]
  Sea $p$ periódico de periodo $n$. Un punto $z$ es \emph{asintótico hacia
  adelante} a $p$ si $\lim_{i \to \infty} f^{in}(z) = p$. El \emph{conjunto
  estable} de $p$, denotado por $W^s(p)$, está formado por todos los puntos
  asintóticos hacia adelante a~$p$.
\end{definicion}

Si $p$ es no periódico, también podríamos definir puntos asintóticos hacia
adelante de la siguiente forma, $\sigma(f^i(z), f^i(p)) \to 0$ cuando $i \to
\infty$. Además, si $f$ tiene inversa, podemos considerar puntos
\emph{asintóticos hacia atrás} haciendo tender $i \to -\infty$ en la definición
de encima. Al conjunto de puntos asintóticos hacia atrás a $p$ se le llama el
conjunto \emph{inestable} de $p$ y se denota por~$W^u(p)$.

\begin{definicion}
  Dada una función $f$ y un punto $z$, la \emph{órbita hacia adelante} de $z$
  es el conjunto de puntos $z$, $f(z)$, $f^2(z)$, $f^3(z)$, $\ldots$ y la
  denotaremos por $\orbit^+(z)$. Si $f$ es un homeomorfismo, la \emph{órbita
  completa} de $z$, $\orbit(z)$, es el conjunto de puntos $f^n(z)$ con $n \in
  \Z$, y la \emph{órbita hacia} atrás de $z$, $\orbit^-(z)$, es el conjunto de
  puntos $z$, $f^{-1}(z)$, $f^{-2}(z)$,~\dots
\end{definicion}

% TODO: In the future when the definition is removed from the appendix and added
% to a section, take into account $\infty$ as a possible $z$. Then f' won't be
% good enough for the definition.
\begin{definicion}[Punto crítico]
  Un punto $z$ es un \emph{punto crítico} de $f$ si $f'(z) = 0$. El punto
  crítico es no degenerado si $f''(z) \neq 0$. El punto crítico es degenerado
  si~$f''(z) = 0$.
\end{definicion}

% --- Ideas descartadas --------------------------------------------------------

% Los puntos fijos tenienen un rol esencial en la dinámica compleja. El siguiente
% resultado nos muestra una de las maneras en la que los puntos fijos son
% esenciales.
%
% \begin{proposicion}
%   Sea $R$ una función racional y $z_0$ un punto en $ExtC$. Entonces, si
%   $(z_n)_{n \in \N}$ es convergente, $\lim_{n \to \infty} z_n = w$ es un punto
%   fijo de $R$.
% \end{proposicion}
% \begin{proof}
%   Como $R$ es continua en su dominio de definición, se cumple lo siguiente:
%   \begin{equation*}
%     w = \lim_{n \to \infty} z_n = \lim_{n \to \infty} R(z_{n-1}) = R \big( \lim_{n \to \infty} z_{n-1} \big) = R(w).
%   \end{equation*}
%   Por tanto, $w = R(w)$, es decir, $w$ es un punto fijo de $R$.
% \end{proof}
